% !TEX encoding = UTF-8 Unicode

\documentclass[final,12pt,a4paper]{article}
\usepackage{wrapfig}
\usepackage[font=footnotesize]{caption}
\usepackage{fullpage}
\usepackage{els}
%\usepackage{cassirer}
\usepackage{ii}

\DeclareSourcemap{
  \maps[datatype=bibtex]{
    \map{
      \step[fieldsource=entrykey, match=DM,fieldset=shorthand, null]
      \step[fieldsource=entrykey, match=ASWB,fieldset=shorthand, null]
      \step[fieldsource=entrykey, match=ESC,fieldset=shorthand, null]
      \step[fieldsource=options,match={skipbib=true},replace={skipbib=false}]
    }
  }
}



\newcommand{\DB}{Dewan-Beran\xspace}
\newcommand{\DBB}{Dewan-Beran-Bell\xspace}
\newcommand{\EP}{\Ehr paradox\xspace}
\newcommand{\AP}{\jt{Annalen der Physik}\xspace}
\renewcommand{\BP}{Bell paradox\xspace}
\newcommand{\TDE}{transverse Doppler effect\xspace}
\newcommand{\TBS}{thread-between-spaceships\xspace}
\newcommand{\RD}{rotating disk\xspace}
\newcommand{\vIgna}{von \Igna\xspace}
\newcommand{\MM}{Michelson-Morley\xspace}
\newcommand{\TE}{thought experiment\xspace}


\hypersetup{draft}


\setlength{\textfloatsep}{10pt plus 1.0pt minus 2.0pt}


\title{Is Time Dilation \s{Real}{}? Einstein and the Transverse Doppler Effect}

\begin{document}
\maketitle


\begin{abstract}
As early as 1907, Einstein realized that the second-order transverse Doppler redshift reported by Stark in the spectral lines of fast-moving ions in canal rays could serve as a direct experimental test of time dilation. He rejected Stark’s \emph{dynamical explanation} of the effect as a real frequency change of the moving ions. Instead, assuming that ions of the same kind emit the same spectral lines when at relative rest, Einstein offered a \emph{kinematical explanation} of the frequency shift as a consequence of time dilation. He labeled the effect \textit{apparent}, since it disappears for the co-moving observer. At the same time, he regarded it as \textit{real}, in the sense that it would not occur if the old kinematics held. From the 1920s onward, Einstein argued that the behavior of atomic clocks calls for a dynamical explanation. The present paper shows, however, that even if an special relativistic explanation of the spectral identity of atoms were available, it would merely reinforce the claim that relativistic time dilation admits a purely kinematical interpretation. It concludes that modern dynamists misread Einstein’s plea for a dynamical account of clock behavior as a demand for \emph{explanation} of the time dilation, whereas it was primarily motivated by the problem of its \emph{confirmation}.
 \end{abstract}

\begin{keywords}
Albert Einstein \sep time dilation \sep special relativity \sep kinematical versus dynamical explanations
\end{keywords}

\vspace{1cm}

\begin{center}
Data availability statement: not applicable.
\end{center}



\thispagestyle{empty}

\newpage

\section*{Ethical Statement}

\begin{itemize}
\item Funding: None
\item Conflict of Interest: None declared
\item  Ethical approval: Not required
\item Informed consent: Not required
\item Author contribution: I am the sole author of this paper

\end{itemize}


\end{document}